
\documentclass{article}
\usepackage{colortbl}
\usepackage{makecell}
\usepackage{multirow}
\usepackage{supertabular}

\begin{document}

\newcounter{utterance}

\centering \large Interaction Transcript for game `wordle\_withclue', experiment `medium\_frequency\_words\_clue\_no\_critic', episode 3 with reflexion{-}llama{-}3.1{-}8b{-}instruct.
\vspace{24pt}

{ \footnotesize  \setcounter{utterance}{1}
\setlength{\tabcolsep}{0pt}
\begin{supertabular}{c@{$\;$}|p{.15\linewidth}@{}p{.15\linewidth}p{.15\linewidth}p{.15\linewidth}p{.15\linewidth}p{.15\linewidth}}
    \# & \multicolumn{2}{c}{Player} && \multicolumn{2}{c}{Game Master} \\
    \hline

    \theutterance \stepcounter{utterance}  
    & & & \multicolumn{4}{p{0.6\linewidth}}{
        \cellcolor[rgb]{0.9,0.9,0.9}{
            \makecell[{{p{\linewidth}}}]{
                \texttt{\tiny{[P1$\langle$GM]}}
                \texttt{You have attempted this same task before and did not succeed. These are the plans you wrote after those attempts.} \\
\texttt{Trial \#1: New plan: The last episode got aborted because "lovely" is not a valid guess since it contains the letter "y" which is red in the guess\_feedback. Moreover, I should have paid closer attention to the clue "sweetheart" and the guess\_feedback to come up with a more informed guess. Next time, I will make sure to produce a five{-}letter word that does not contain red letters (o, v, l, y) and fits the theme of affection or love. I will also carefully consider the clue and the guess\_feedback to come up with a more accurate guess. For example, I could try a word like "heart" or "hearts".} \\
\\ 
\texttt{Trial \#2: New plan: The last episode got aborted because "dear" is not a valid guess since it contains only 4 letters, not 5. Moreover, I should have paid closer attention to the clue "sweetheart" and the guess\_feedback to come up with a more informed guess. Next time, I will make sure to produce a five{-}letter word that does not contain red letters (o, v, l, y) and fits the theme of affection or love. I will also carefully consider the clue and the guess\_feedback to come up with a more accurate guess. For example, I could try a word like "hearts" or "honey".} \\
\\ 
\texttt{Trial \#3: New plan: The last two episodes got aborted because I failed to produce a valid five{-}letter word that fits the theme of affection or love. In both attempts, I made the mistake of either including a red letter or producing a word with fewer than five letters. Moreover, I should have paid closer attention to the clue "sweetheart" and the guess\_feedback to come up with a more informed guess. Next time, I will make sure to produce a five{-}letter word that does not contain red letters (o, v, l, y) and carefully consider the clue and the guess\_feedback to come up with a more accurate guess. I will also double{-}check the length of my word to ensure it has exactly five letters. For example, I could try a word like "honey" or "loves".} \\
\\ 
\\ 
\texttt{You are a language wizard who likes to guess words by using the given rules.} \\
\\ 
\texttt{Welcome to Wordle! You have six attempts to guess the target word, a valid English word of five lowercase letters (a{-}z). Please use the tags "explanation:" and "guess:" to provide a concise explanation for each guess.} \\
\\ 
\texttt{To help you make an informed guess, you will receive a clue for the word, such as} \\
\texttt{clue: snowy white.} \\
\\ 
\texttt{Here is an example guess based on the clue:} \\
\texttt{explanation: In the fairy tail Snow White, the girl is killed because she eats a poisoned apple. And the word apple has 5 letters.} \\
\texttt{guess: apple} \\
\\ 
\texttt{After each guess, your answer will be validated, and you will receive feedback indicating which letters are correct (green), which letters are correct but in the wrong position (yellow), and which letters are incorrect (red). This feedback can be useful in determining which letters to include or exclude in your next guess.} \\
\\ 
\texttt{For example, the feedback for "apple" might be:} \\
\texttt{guess\_feedback: a<yellow> p<yellow> p<green> l<yellow> e<red>} \\
\\ 
\texttt{The explanation should contain details about how both the guess\_feedback and clue work together in tandem to formulate a new guess.} \\
\\ 
\texttt{Do not generate any other text.} \\
\\ 
\texttt{Let's begin with your first guess.} \\
\\ 
\texttt{clue: sweetheart} \\
            }
        }
    }
     \\ \\

    \theutterance \stepcounter{utterance}  
    & \multicolumn{4}{p{0.6\linewidth}}{
        \cellcolor[rgb]{0.561,0.737,0.561}{
            \makecell[{{p{\linewidth}}}]{
                \texttt{\tiny{[P1$\rangle$GM]}}
                \texttt{explanation: Based on the clue "sweetheart," I will try to come up with a five{-}letter word that fits the theme of affection or love. Since the clue is related to a term of endearment, I will focus on words that convey a sense of closeness or affection. I will also pay close attention to the guess\_feedback to adjust my next guess accordingly.} \\
\\ 
\texttt{guess: honey} \\
            }
        }
    }
    & & \\ \\

    \theutterance \stepcounter{utterance}  
    & & & \multicolumn{2}{p{0.3\linewidth}}{
        \cellcolor[rgb]{0.9,0.9,0.9}{
            \makecell[{{p{\linewidth}}}]{
                \texttt{\tiny{[GM$|$GM]}}
                \texttt{attempts = 1} \\
\texttt{target\_word = flame} \\
\texttt{target\_word\_clue = sweetheart} \\
\texttt{guess = honey} \\
\texttt{guess\_feedback = h<red> o<red> n<red> e<yellow> y<red>} \\
            }
        }
    }
    & & \\ \\

    \theutterance \stepcounter{utterance}  
    & & & \multicolumn{4}{p{0.6\linewidth}}{
        \cellcolor[rgb]{0.9,0.9,0.9}{
            \makecell[{{p{\linewidth}}}]{
                \texttt{\tiny{[P1$\langle$GM]}}
                \texttt{guess\_feedback: h<red> o<red> n<red> e<yellow> y<red>} \\
\\ 
\texttt{Provide your response only in this format.} \\
\texttt{explanation: details} \\
\texttt{guess: word} \\
            }
        }
    }
     \\ \\

    \theutterance \stepcounter{utterance}  
    & \multicolumn{4}{p{0.6\linewidth}}{
        \cellcolor[rgb]{0.561,0.737,0.561}{
            \makecell[{{p{\linewidth}}}]{
                \texttt{\tiny{[P1$\rangle$GM]}}
                \texttt{explanation: Since the guess "honey" received mostly red letters, I will focus on the yellow letter "e" which is present in the target word. I will also avoid the letters "h", "o", "n", and "y" which are red in the guess\_feedback. Considering the clue "sweetheart," I will try to come up with a new five{-}letter word that fits the theme of affection or love and does not contain the red letters.} \\
\\ 
\texttt{guess: hearts} \\
            }
        }
    }
    & & \\ \\

    \theutterance \stepcounter{utterance}  
    & & & \multicolumn{2}{p{0.3\linewidth}}{
        \cellcolor[rgb]{0.9,0.9,0.9}{
            \makecell[{{p{\linewidth}}}]{
                \texttt{\tiny{[GM$|$GM]}}
                \texttt{Error: The length of the guessed word is not 5.} \\
            }
        }
    }
    & & \\ \\

    \theutterance \stepcounter{utterance}  
    & & & \multicolumn{2}{p{0.3\linewidth}}{
        \cellcolor[rgb]{0.9,0.9,0.9}{
            \makecell[{{p{\linewidth}}}]{
                \texttt{\tiny{[GM$|$GM]}}
                \texttt{game\_result = ABORT} \\
            }
        }
    }
    & & \\ \\

    \theutterance \stepcounter{utterance}  
    & & \multicolumn{4}{p{0.6\linewidth}}{
        \cellcolor[rgb]{0.95,0.95,0.95}{
            \makecell[{{p{\linewidth}}}]{
                \texttt{\tiny{[P1$\leftrightarrow$P]}}
                \texttt{Post{-}episode Reflexion (not sent to GM):} \\
\texttt{New plan: The last episode got aborted because "hearts" is not a valid guess since it contains the letter "s" which is not present in the clue "sweetheart". Moreover, I should have paid closer attention to the clue "sweetheart" and the guess\_feedback to come up with a more informed guess. Next time, I will make sure to produce a five{-}letter word that does not contain red letters (o, v, l, y) and carefully consider the clue and the guess\_feedback to come up with a more accurate guess. I will also double{-}check the presence of letters in the clue and the guess\_feedback to ensure that my guess is more informed. For example, I could try a word like "honey" or "loves" but with a more careful consideration of the clue and the guess\_feedback.} \\
            }
        }
    }
    & \\ \\

\end{supertabular}
}

\end{document}
